深度学习与自然语言处理核心概念综合梳理
本文档系统梳理了从统计语言模型到现代大模型（Transformer）的完整技术栈，涵盖数据表示、经典架构、数学优化、工程实践与现代注意力机制。

第一部分：序列与输入表示（Tokenization & Embeddings）
1. 分词（Tokenization）与子词模型（Subword Models）
目标：将原始文本转化为模型可处理的基本单元（Token），解决词表过大（OOV）及序列过长问题。

BPE（字节对编码）：基于相邻共现频率进行迭代合并。

初始化：将单词拆分为字符级单元（如 
h
u
g
→
h
,
u
,
g
hug→h,u,g）。

统计频率：计算语料中所有相邻符号对 
(
x
,
y
)
(x,y) 的出现次数。

合并规则：将频率最高的对合并为新子词（如 
h
+
u
→
h
u
h+u→hu），重复直至词表达到预设大小。

数学本质：通过贪心策略最大化语料的压缩率。

Unigram 语言模型：基于概率统计进行剪枝。

初始化一个极大的可能子词词表。

通过 EM 算法估计每个子词的出现概率 
P
(
x
)
P(x)。

衡量移除某子词对整体似然（Likelihood）的损失，优先剔除损失增量最小（即频率最低）的子词。

2. 词元嵌入（Token Embeddings）
目标：将离散的 Token ID 映射为可学习的连续稠密向量，以便进行数学运算。

数学表示：设词表大小为 
∣
V
∣
∣V∣，嵌入维度为 
d
d。嵌入矩阵为 
E
∈
R
∣
V
∣
×
d
E∈R 
∣V∣×d
 。

每个 Token 
w
i
w 
i
​
  对应矩阵中的一行向量 
v
i
=
E
[
i
,
:
]
v 
i
​
 =E[i,:]。

“可学习（Learned）”的数学含义：在反向传播中，嵌入矩阵 
E
E 作为参数参与梯度更新。

根据链式法则，损失 
L
L 对嵌入向量的偏导为 
∂
L
∂
v
i
∂v 
i
​
 
∂L
​
 ，驱动向量在空间中位移，使得语义相近的词在几何距离上靠近。

第二部分：经典语言模型（从统计到神经）
1. Ngram 模型（Bigram / Trigram）
核心假设：马尔可夫假设（Markov Assumption），当前词仅依赖于前 
N
−
1
N−1 个词。

Bigram 数学定义（
N
=
2
N=2）：
P
(
w
t
∣
w
1
:
t
−
1
)
≈
P
(
w
t
∣
w
t
−
1
)
=
Count
(
w
t
−
1
,
w
t
)
Count
(
w
t
−
1
)
P(w 
t
​
 ∣w 
1:t−1
​
 )≈P(w 
t
​
 ∣w 
t−1
​
 )= 
Count(w 
t−1
​
 )
Count(w 
t−1
​
 ,w 
t
​
 )
​
 

最大似然估计（MLE）：通过数频次计算条件概率，选择使得训练数据出现概率最大的参数集。

对数空间（Log-Space）：

联合概率连乘 
P
=
∏
p
i
P=∏p 
i
​
  极易导致数值下溢。利用 
log
⁡
(
a
b
)
=
log
⁡
a
+
log
⁡
b
log(ab)=loga+logb 转换。

目标：最大化对数似然 
∑
log
⁡
P
(
w
i
∣
w
i
−
1
)
∑logP(w 
i
​
 ∣w 
i−1
​
 )。

困惑度（Perplexity, PPL）：评估模型质量的指标。
P
P
L
=
exp
⁡
(
−
1
N
∑
i
=
1
N
log
⁡
P
(
w
i
∣
w
i
−
1
)
)
PPL=exp(− 
N
1
​
 ∑ 
i=1
N
​
 logP(w 
i
​
 ∣w 
i−1
​
 ))

物理意义：模型预测下一个词时，平均有多少个“等概率”候选词在竞争（值越低越好）。

局限性：有限上下文窗口（Fixed Context），无法捕捉长距离依赖，且数据极度稀疏（Sparsity）。

2. 前馈神经语言模型（NNLM）与 Deep CBoW
革命点：引入神经网络和词嵌入，通过参数共享解决稀疏性问题。

前馈 NNLM：将前 
N
−
1
N−1 个词的嵌入向量拼接（Concat），经过非线性隐藏层（
t
a
n
h
tanh）后预测下一个词。

优势：保留了词序信息。

Deep CBoW（深度连续词袋模型）：

核心操作：将上下文词向量进行平均/求和（Pooling），彻底丢弃词序信息。

数学表达：
设上下文词向量为 
v
c
1
,
.
.
.
,
v
c
k
v 
c1
​
 ,...,v 
ck
​
 ，池化后得到隐含表示：
h
=
1
k
∑
i
=
1
k
v
c
i
h= 
k
1
​
 ∑ 
i=1
k
​
 v 
ci
​
 

“深度（Deep）”体现：在 
h
h 与输出 Softmax 之间插入多个全连接隐藏层（如 
h
1
=
ReLU
(
W
1
h
+
b
1
)
h 
1
​
 =ReLU(W 
1
​
 h+b 
1
​
 )），引入非线性特征交叉。

本质：词袋特征 + 多层非线性分类器。

第三部分：训练机制与数学优化（如何更新参数）
1. 训练三步走（Loss, Graph, Gradients）
目标：通过反向传播（Backpropagation）更新权重，最小化损失。

选择损失函数（Cross-Entropy Loss）：衡量预测分布 
q
(
x
)
q(x) 与真实分布 
p
(
x
)
p(x) 的差距。
L
=
−
∑
i
p
(
x
i
)
log
⁡
q
(
x
i
)
L=−∑ 
i
​
 p(x 
i
​
 )logq(x 
i
​
 )
对于分类任务，真实标签为 One-hot，上式简化为 
L
=
−
log
⁡
q
(
x
t
a
r
g
e
t
)
L=−logq(x 
target
​
 )。

构建可微分计算图（Differentiable Graph）：框架将矩阵乘法、激活函数（ReLU, Softmax）等算子记录为节点，张量流为边。

求梯度（Chain Rule）：从损失节点逆向传播。
∂
L
∂
w
=
∂
L
∂
y
^
⋅
∂
y
^
∂
z
⋅
∂
z
∂
w
∂w
∂L
​
 = 
∂ 
y
^
​
 
∂L
​
 ⋅ 
∂z
∂ 
y
^
​
 
​
 ⋅ 
∂w
∂z
​
 

其中 
z
=
W
x
+
b
z=Wx+b，通过链式法则将误差逐层分配给所有参数矩阵。

2. 梯度消失与门控解决方案
问题：标准 RNN 中，梯度 
∂
h
t
∂
h
t
−
1
=
W
T
⋅
diag
(
tanh
⁡
′
)
∂h 
t−1
​
 
∂h 
t
​
 
​
 =W 
T
 ⋅diag(tanh 
′
 ) 若小于 1 则连乘导致梯度消失，无法建模长序列。

解决方案 1：加法连接（Residual Connection）
Output
=
F
(
x
)
+
x
Output=F(x)+x

反向传播：
∂
Output
∂
x
=
∂
F
∂
x
+
1
∂x
∂Output
​
 = 
∂x
∂F
​
 +1

数学优势：即使 
F
F 层的梯度消失，常数项 
1
1 仍能保证梯度无损回传。

解决方案 2：门控机制（Gating）—— LSTM / GRU

第四部分：门控循环网络详解（LSTM 与 GRU）
1. LSTM（长短期记忆网络，1997）
核心思想：引入独立的细胞状态（Cell State, 
C
t
C 
t
​
 ）作为“传送带”，通过门控控制信息增删。

遗忘门（Forget Gate, 
f
t
f 
t
​
 ）：决定丢弃多少旧记忆。
f
t
=
σ
(
W
f
⋅
[
h
t
−
1
,
x
t
]
+
b
f
)
f 
t
​
 =σ(W 
f
​
 ⋅[h 
t−1
​
 ,x 
t
​
 ]+b 
f
​
 )

当 
f
t
≈
1
f 
t
​
 ≈1 时，旧状态 
C
t
−
1
C 
t−1
​
  完全保留，实现长期记忆。

输入门（Input Gate, 
i
t
i 
t
​
 ）与候选细胞（
C
~
t
C
~
  
t
​
 ）：决定写入多少新信息。
i
t
=
σ
(
W
i
⋅
[
h
t
−
1
,
x
t
]
+
b
i
)
i 
t
​
 =σ(W 
i
​
 ⋅[h 
t−1
​
 ,x 
t
​
 ]+b 
i
​
 )
C
~
t
=
tanh
⁡
(
W
C
⋅
[
h
t
−
1
,
x
t
]
+
b
C
)
C
~
  
t
​
 =tanh(W 
C
​
 ⋅[h 
t−1
​
 ,x 
t
​
 ]+b 
C
​
 )

细胞状态更新（核心直通路径）：
C
t
=
f
t
⊙
C
t
−
1
+
i
t
⊙
C
~
t
C 
t
​
 =f 
t
​
 ⊙C 
t−1
​
 +i 
t
​
 ⊙ 
C
~
  
t
​
 

数学解析：若 
f
t
=
1
f 
t
​
 =1 且 
i
t
=
0
i 
t
​
 =0，则 
C
t
=
C
t
−
1
C 
t
​
 =C 
t−1
​
 ，梯度在此路径上的导数为 
1
1，完美解决了梯度消失。

输出门（Output Gate, 
o
t
o 
t
​
 ）：决定对外暴露多少细胞状态信息。
h
t
=
o
t
⊙
tanh
⁡
(
C
t
)
h 
t
​
 =o 
t
​
 ⊙tanh(C 
t
​
 )

2. GRU（门控循环单元，2014）
简化版：将 LSTM 的 3 门结构压缩为 2 门（更新门 
z
t
z 
t
​
  和重置门 
r
t
r 
t
​
 ），无独立细胞状态。

更新门（
z
t
z 
t
​
 ）：控制历史状态保留比例（相当于 LSTM 的遗忘门 + 输入门合并）。

重置门（
r
t
r 
t
​
 ）：控制忽略历史状态的程度。

最终更新：
h
t
=
z
t
⊙
h
t
−
1
+
(
1
−
z
t
)
⊙
h
~
t
h 
t
​
 =z 
t
​
 ⊙h 
t−1
​
 +(1−z 
t
​
 )⊙ 
h
~
  
t
​
 。

同样通过 
z
t
≈
1
z 
t
​
 ≈1 实现记忆直通。

第五部分：现代工程实践（Practical Considerations）
概念	数学/技术原理	作用
权重初始化（He Initialization）	方差设为 
2
/
fan_in
2/fan_in（适配 ReLU）	防止前向/后向传播中信号方差指数级衰减或爆炸
优化器（AdamW）	带动量的自适应学习率：
m
t
=
β
1
m
t
−
1
+
(
1
−
β
1
)
g
t
m 
t
​
 =β 
1
​
 m 
t−1
​
 +(1−β 
1
​
 )g 
t
​
 ；
v
t
=
β
2
v
t
−
1
+
(
1
−
β
2
)
g
t
2
v 
t
​
 =β 
2
​
 v 
t−1
​
 +(1−β 
2
​
 )g 
t
2
​
 ；参数更新：
θ
t
+
1
=
θ
t
−
η
m
t
v
t
+
ϵ
θ 
t+1
​
 =θ 
t
​
 −η 
v 
t
​
 
​
 +ϵ
m 
t
​
 
​
 	加速收敛，跨越鞍点，解耦权重衰减
学习率调度（Warmup + Cosine）	预热：
lr
=
step
/
warmup_steps
×
peak
lr=step/warmup_steps×peak；余弦退火：
lr
=
peak
⋅
0.5
(
1
+
cos
⁡
(
π
t
T
)
)
lr=peak⋅0.5(1+cos(π 
T
t
​
 ))	前期稳定梯度统计，后期精细微调
过拟合对抗（Dropout）	训练时以概率 
p
p 将神经元输出置零：
y
=
mask
⊙
x
y=mask⊙x	强迫网络学习冗余表示，增强泛化性
梯度累积	多批次（Micro-batches）计算梯度求和，更新前统一除以累积步数 
N
N：
g
=
1
N
∑
g
i
g= 
N
1
​
 ∑g 
i
​
 	用小显存模拟大 Batch Size 的效果
第六部分：序列建模的三大范式
1. 循环（Recurrence, RNN/LSTM）
公式：
h
t
=
f
(
h
t
−
1
,
x
t
)
h 
t
​
 =f(h 
t−1
​
 ,x 
t
​
 )

特性：串行计算，理论感受野无限，但无法并行且存在长距离遗忘风险。

2. 卷积（Convolution, TCN/WaveNet）
公式：
y
i
=
∑
j
=
0
k
−
1
W
j
⋅
x
i
+
j
y 
i
​
 =∑ 
j=0
k−1
​
 W 
j
​
 ⋅x 
i+j
​
 （卷积核大小为 
k
k）

特性：完全并行计算，梯度稳定，但感受野受限，需堆叠多层或使用空洞卷积。

3. 注意力（Attention, Transformer）
公式（缩放点积）：
Attention
(
Q
,
K
,
V
)
=
softmax
(
Q
K
T
d
k
)
V
Attention(Q,K,V)=softmax( 
d 
k
​
 
​
 
QK 
T
 
​
 )V

特性：全局感受野（
O
(
1
)
O(1) 距离），并行计算，但计算复杂度为 
O
(
n
2
)
O(n 
2
 )（显存瓶颈）。

第七部分：Transformer 注意力机制数学与维度分析
1. 缩放点积注意力（Scaled Dot-Product Attention）
输入：Query 
Q
Q，Key 
K
K，Value 
V
V。

步骤：

计算得分（Score）：
S
=
Q
⋅
K
T
S=Q⋅K 
T
 。
缩放（Scale）：除以 
d
k
d 
k
​
 
​
  防止点积过大导致 Softmax 梯度饱和。
掩码（Mask，可选）：将无效位置设为 
−
∞
−∞（
S
=
S
.
masked_fill
(
m
a
s
k
=
0
,
v
a
l
u
e
=
−
1
e
9
)
S=S.masked_fill(mask=0,value=−1e9)）。
Softmax 归一化：
A
=
softmax
(
S
)
A=softmax(S)，得到权重矩阵。
加权求和：
Context
=
A
⋅
V
Context=A⋅V。
2. 多头注意力（Multi-Head Attention）
数学原理：

设 
d
model
d 
model
​
  为总维度，
h
h 为头数，每个头的维度 
d
k
=
d
model
/
h
d 
k
​
 =d 
model
​
 /h。

对输入 
X
X 进行 
h
h 次并行投影（线性层 
W
i
Q
,
W
i
K
,
W
i
V
W 
i
Q
​
 ,W 
i
K
​
 ,W 
i
V
​
 ）：
head
i
=
Attention
(
X
W
i
Q
,
X
W
i
K
,
X
W
i
V
)
head 
i
​
 =Attention(XW 
i
Q
​
 ,XW 
i
K
​
 ,XW 
i
V
​
 )

拼接所有头并经过输出投影 
W
O
W 
O
 ：
MultiHead
(
X
)
=
Concat
(
head
1
,
.
.
.
,
head
h
)
W
O
MultiHead(X)=Concat(head 
1
​
 ,...,head 
h
​
 )W 
O
 

维度流（以 
X
∈
R
b
a
t
c
h
×
s
e
q
×
d
m
o
d
e
l
X∈R 
batch×seq×d 
model
​
 
  为例）：

线性层后：
Q
,
K
,
V
∈
R
b
a
t
c
h
×
s
e
q
×
d
m
o
d
e
l
Q,K,V∈R 
batch×seq×d 
model
​
 
 。

拆分多头并转置：
Q
,
K
,
V
∈
R
b
a
t
c
h
×
h
e
a
d
s
×
s
e
q
×
d
k
Q,K,V∈R 
batch×heads×seq×d 
k
​
 
 。

Attention 计算后：
Context
∈
R
b
a
t
c
h
×
h
e
a
d
s
×
s
e
q
×
d
k
Context∈R 
batch×heads×seq×d 
k
​
 
 。

转置并合并：
Concat
∈
R
b
a
t
c
h
×
s
e
q
×
d
m
o
d
e
l
Concat∈R 
batch×seq×d 
model
​
 
 。

3. 代码中的工程对齐（Dropout 与 Mask）
Dropout 位置：在 
Softmax
Softmax 计算出的注意力权重 
A
A 后应用，用于正则化注意力分布，防止过拟合。

因果掩码（Causal Mask）：在自回归生成（GPT）中，通过生成上三角矩阵 
M
M（
M
i
j
=
0
M 
ij
​
 =0 if 
i
≥
j
i≥j else 
−
∞
−∞）屏蔽未来信息，确保位置 
t
t 只能看到 
1
1 至 
t
t 的 Token。

总结：技术演进逻辑图
统计计数 (Ngram)
→
稀疏性
神经嵌入 (NNLM/CBoW)
→
固定窗口/串行
门控 (LSTM/GRU)
统计计数 (Ngram) 
稀疏性
​
 神经嵌入 (NNLM/CBoW) 
固定窗口/串行
​
 门控 (LSTM/GRU)

→
串行瓶颈/无法并行
注意力 (Transformer)
→
平方复杂度
2026前沿 (Mamba/稀疏注意力)
串行瓶颈/无法并行
​
 注意力 (Transformer) 
平方复杂度
​
 2026前沿 (Mamba/稀疏注意力)

核心思想闭环：所有的进步（从门控到残差，从LSTM到Attention）本质上都是在如何更高效地传递梯度（解决梯度消失）和如何更广泛地聚合上下文（打破有限窗口与串行依赖）之间寻找平衡。